\documentclass[11pt]{article}
\usepackage[T1]{fontenc}
\usepackage[margin=27mm]{geometry}
\usepackage{amsmath,amssymb,amsthm,mathtools}
\usepackage{hyperref}
\newtheorem{theorem}{Theorem}
\newtheorem{lemma}{Lemma}
\newtheorem{proposition}{Proposition}
\newtheorem{corollary}{Corollary}
\newcommand{\su}{\mathfrak{su}}
\newcommand{\tr}{\operatorname{tr}}
\newcommand{\diag}{\operatorname{diag}}
\title{The exact two-cross scalar--Gauss domain is nonempty at every positive coupling}
\author{Independent exact derivation}
\date{5 September 2026}
\begin{document}
\maketitle

\section{Literal data, residuals, and coordinate maps}

All complex vector spaces below retain their underlying real structures.
The adjoint is denoted by $*$, and $[X,Y]=XY-YX$. Keep exactly
\[
 H_0=\begin{pmatrix}0&0&0\\0&-1/2&1/(2\sqrt2)\\
 0&1/(2\sqrt2)&1/2\end{pmatrix},\qquad
 q=\frac1{2\sqrt2}\begin{pmatrix}1\\0\\0\end{pmatrix},\qquad
 A_0=\sqrt2\begin{pmatrix}0&-1&0\\1&0&0\\0&0&0\end{pmatrix}.
\]
The original constants are $m^2=\Gamma=1$ and $\kappa=6$.
For $Y_{\rm tf}=Y-(\tr Y)I_3/3$, define
\[
 \mathcal I_HY=([H,[H,Y]]+qq^*Y+Yqq^*)_{\rm tf},\qquad
 \mathcal K_{Z_1,Z_2}Y=-\sum_{j=1}^2[Z_j,[Z_j,Y]].
\]
Direct multiplication gives
\[
 F_0=[A_0,[A_0,H_0]]=-2H_0+\diag(-2,1,1),\quad
 H_0^2=\frac38\diag(0,1,1),\quad q^*q=\frac18.
\]
The scalar--Gauss domain in question consists of
\[
 (t,\omega,C_0,Z_{1,0},Z_{2,0}),\quad t=g^2>0,\quad
 \omega\in\mathbb R\setminus\{0\},\quad C_0,Z_{j,0}\in\su(3),
\]
where $P=qq^*/(q^*q)=\diag(1,0,0)$ and
\[
 PZ_{j,0}P=(I_3-P)Z_{j,0}(I_3-P)=0,
\]
satisfying exactly the following three full residual equations:
\begin{align}
 [C_0,[C_0,H_0]]-\sum_j[Z_{j,0},[Z_{j,0},H_0]]-F_0&=0,
 \label{sgd:literal-H}\\
 (C_0^2-Z_{1,0}^2-Z_{2,0}^2)q+2q&=0,
 \label{sgd:literal-q}\\
 (t\mathcal I_{H_0}+2\mathcal K_{Z_{1,0},Z_{2,0}})C_0
             -4\omega[Z_{1,0},Z_{2,0}]&=0.
 \label{sgd:literal-G}
\end{align}
Equation~\eqref{sgd:literal-G} is the original Gauss equation multiplied
by its retained positive $t$; division by the same $t$ is its inverse.
Neither spatial equation is imposed anywhere in the construction.

The exact maps to and from the classified frame are as follows. Retain
\[
 r=\frac{\sqrt6}{4},\quad \delta=\frac1{2r}=\sqrt{\frac23},\quad
 x=\frac12,\quad d_0=\frac1{2\sqrt2},\quad
 S=\frac1{\sqrt{2r(r+x)}}
       \begin{pmatrix}d_0&-(r+x)\\r+x&d_0\end{pmatrix},\quad
 U=\diag(1,S).
\]
The identity $r^2=x^2+d_0^2$ gives $S^*S=I_2$, $\det S=1$;
multiplication of $H_0$ by its two displayed lower columns gives
\[
 U^*H_0U=\diag(0,r,-r),\qquad U^*q=q.
\]
Choose any $\zeta\in\mathbb C$ with $|\zeta|=1$, and set
\[
 \gamma=-\zeta^{-2},\qquad
 W=\frac\zeta{\sqrt2}\begin{pmatrix}\gamma&\gamma\\1&-1\end{pmatrix}
  =\frac1{\sqrt2}\begin{pmatrix}-\zeta^{-1}&-\zeta^{-1}\\
                         \zeta&-\zeta\end{pmatrix},\qquad
 V=\diag(1,W),\quad B=UV.
\]
Here $|\gamma|=1$ and $\zeta^2=-\overline\gamma$, retaining exactly
the original phase relation. Every unit $\gamma$ occurs, with its two
possible square roots $\zeta$. The displayed entries give
\[
 W^*W=I_2,\quad \det W=1,\quad
 W^*\diag(r,-r)W=\begin{pmatrix}0&r\\r&0\end{pmatrix}.
\]
Thus, simultaneously for all data,
\begin{equation}
 Y'=B^*Y_0B,\qquad Y_0=BY'B^*,\qquad
 H'=\begin{pmatrix}0&0&0\\0&0&r\\0&r&0\end{pmatrix},\quad
 q'=q,\quad F'=-2H'+\diag(-2,1,1).
 \label{sgd:frame}
\end{equation}
These are inverse maps, with no constant or unknown rescaled. In fact
$B^*(XY)B=(B^*XB)(B^*YB)$, the same equality holds for brackets,
trace and adjoints, and $Bq=q$. Therefore $\mathcal I$ and $\mathcal K$
and every full residual in \eqref{sgd:literal-H}--\eqref{sgd:literal-G}
are conjugated exactly (the vector residual is multiplied by $B^*$).
This proves that zero residuals in the classified frame give zero residuals
in the literal frame, with the displayed inverse recovering every datum.
The block-diagonal $B$ commutes with $P$, preserving the two-cross class.

\section{Every scalar and Gauss entry in the classified frame}

For $v,w\in\mathbb C^2$ write
\[
 X(v)=\begin{pmatrix}0&-v^*\\v&0_{2\times2}\end{pmatrix},\qquad
 Z_1=X(v),\quad Z_2=X(w),\quad
 C=ie\diag(1,0,-1),\quad e\in\mathbb R\setminus\{0\}.
\]
The inverse of $v\mapsto X(v)$ is the lower first column; every
$X(v)$ is traceless anti-Hermitian. In this section all unprimed matrices
are these classified-frame matrices, with literal inverse~\eqref{sgd:frame}.
Define
\[
 M=vv^*+ww^*=\begin{pmatrix}u&m\\\bar m&v_0\end{pmatrix},\quad
 N=i(vw^*-wv^*),\quad b=\operatorname{Im}(v^*w),\quad
 Q=\diag(1,2),\quad E=\diag(2/3,-5/6).
\]
The subscript on the real scalar $v_0$ distinguishes it from the vector $v$.

Here are the complete residuals, rather than only contractions of them:
\begin{align}
 Q_H&=\begin{pmatrix}
 2-4r\operatorname{Re}m&0&0\\
 0&2r\operatorname{Re}m-1&r(\tr M-e^2+2)\\
 0&r(\tr M-e^2+2)&2r\operatorname{Re}m-1
 \end{pmatrix},\label{sgd:scalar-H}\\
 Q_q&=(\tr M-e^2+2)q,\label{sgd:scalar-q}\\
 Q_G&=i\diag(\gamma_0,\mathcal G),\label{sgd:Gauss-full}\\
 \gamma_0&=et/6+4e(u+2v_0)+8\omega b,\qquad
 \mathcal G=etE-2e(QM+MQ)-4\omega N.\label{sgd:Gauss-entries}
\end{align}
In \eqref{sgd:Gauss-full} the second diagonal block has size two.
To derive every displayed entry, multiplication gives
\[
 X(v)^2=\diag(-v^*v,-vv^*),\quad
 [C,X(v)]=X(-ieQv),\quad
 [X(v),X(w)]=\diag(-2ib,iN).
\]
The first two identities and $H'_{23}=H'_{32}=r$ give
\eqref{sgd:scalar-H} and \eqref{sgd:scalar-q}; all their omitted entries
are zero. The two contributions to $\mathcal I_{H'}C$ are
\[
 ie\diag(0,3/4,-3/4),\qquad
 ie\diag(1/6,-1/12,-1/12),
\]
so $\mathcal I_{H'}C=ie\diag(1/6,2/3,-5/6)$. Also
\[
 \mathcal KC=\diag\bigl(2ie\tr(QM),-ie(QM+MQ)\bigr).
\]
Substitution of these full blocks gives \eqref{sgd:Gauss-full}.
Finally $\tr N=2b$, since $\tr(vw^*)=w^*v=\overline{v^*w}$.
Thus $\tr\mathcal G=-\gamma_0$; the upper Gauss entry vanishes
whenever the full lower block does, and this implication is an identity,
not an omitted equation. The two scalar equations are equivalent to
\begin{equation}
 e^2=2+u+v_0,\qquad \operatorname{Re}m=\delta.
 \label{sgd:scalar-conditions}
\end{equation}

\section{The exact realization of the two Hermitian matrices}

\begin{lemma}\label{sgd:realization-necessity}
For every pair $v,w$, let $D=(v\ w)$ and
$J=\left(\begin{smallmatrix}0&i\\-i&0\end{smallmatrix}\right)$.
Then $M=DD^*$, $N=DJD^*$ and
\[
 M+N=(v-iw)(v-iw)^*,\qquad
 M-N=(v+iw)(v+iw)^*.
\]
If $M$ is positive definite, then
\begin{equation}
 \tr(M^{-1}N)=0,\qquad \det N=-\det M.
 \label{sgd:trace-det}
\end{equation}
\end{lemma}
\begin{proof}
Multiplication of the displayed two-by-two blocks gives all three
factorizations. If $M>0$, $D$ is invertible because
$z^*Mz=\|D^*z\|^2$ has zero kernel. Consequently
$M^{-1}N=D^{-*}JD^*$ is similar to $J$, whose trace is zero
and determinant is $-1$. Taking the trace and determinant proves
\eqref{sgd:trace-det}, with every factor retained.
\end{proof}

\begin{lemma}\label{sgd:rank-one-impossible}
Every classified scalar--Gauss solution has $M>0$.
\end{lemma}
\begin{proof}
The scalar equation gives $m=\delta+iy\ne0$, so $u,v_0>0$
and $M\ne0$. It is already positive semidefinite as $DD^*$.
If $M$ had rank one, the range of $N=DJD^*$ would be contained
in its one-dimensional range. A Hermitian matrix with that range
and its orthogonal complement in its kernel is a real scalar multiple
of $M$, so $N=\nu M$ for some real $\nu$.
Set $k=e/\omega\ne0$. The full Gauss equation reads
\begin{equation}
 N=kT,\qquad
 T=\frac t4 E-\frac12(QM+MQ)
  =\begin{pmatrix}t/6-u&-3m/2\\-3\bar m/2&-5t/24-2v_0\end{pmatrix}.
 \label{sgd:T}
\end{equation}
Its off-diagonal entry gives $\nu=-3k/2$ because $m\ne0$.
The $(1,1)$ entry then gives $t/6-u=-3u/2$, or $t/6=-u/2$,
which contradicts $t>0$ and $u>0$. Thus $M$ has rank two and is
positive definite.
\end{proof}

\begin{theorem}[Complete classified-frame parameterization]
\label{sgd:parametrization}
Define the open real parameter set
\[
 \mathcal D=\{(u,v_0,y)\in\mathbb R^3:
      u>0,\quad d=uv_0-\delta^2-y^2>0,\quad c=4v_0-5u>0\}.
\]
For $(u,v_0,y)\in\mathcal D$, retain
\begin{align}
 m&=\delta+iy,&
 M&=\begin{pmatrix}u&m\\\bar m&v_0\end{pmatrix},&
 t&=\frac{72d}{c},\label{sgd:parameters1}\\
 T&=\frac t4E-\frac12(QM+MQ),&
 k&=\sigma\sqrt{\frac d{-\det T}},&
 e&=\epsilon\sqrt{2+u+v_0},\label{sgd:parameters2}\\
 \omega&=e/k,&N&=kT,& \epsilon,\sigma&\in\{1,-1\}.
 \label{sgd:parameters3}
\end{align}
The denominator $-\det T$ is strictly positive. The two matrices
$A_+=M+N$ and $A_-=M-N$ are positive semidefinite of rank one,
and their $(1,1)$ entries $h_+,h_-$ are strictly positive. Define
\begin{equation}
 a_\pm=\begin{pmatrix}\sqrt{h_\pm}\\(A_\pm)_{21}/\sqrt{h_\pm}\end{pmatrix},
 \quad p_\pm\in\mathbb C,\quad |p_\pm|=1,\quad
 v=\frac{p_+a_++p_-a_-}{2},\quad
 w=\frac{i(p_+a_+-p_-a_-)}{2}.
 \label{sgd:vectors}
\end{equation}
These formulas give exactly all classified-frame solutions of
\eqref{sgd:scalar-H}--\eqref{sgd:Gauss-entries}. In fact they are a
bijection from $\mathcal D\times\{1,-1\}^2\times U(1)^2$ to
that solution set with coordinates $(t,\omega,e,v,w)$.
\end{theorem}
\begin{proof}
First consider any solution. Lemma~\ref{sgd:rank-one-impossible}
gives $M>0$ and $d>0$. Inverting its two-by-two matrix gives
\[
 M^{-1}=\frac1d\begin{pmatrix}v_0&-m\\-\bar m&u\end{pmatrix}.
\]
The trace condition in \eqref{sgd:trace-det}, divided by the retained
nonzero $k$, now gives exactly
\begin{align}
 0=d\tr(M^{-1}T)
 &=v_0(t/6-u)+u(-5t/24-2v_0)+3|m|^2\notag\\
 &=\frac{t(4v_0-5u)}{24}-3d.\label{sgd:trace-condition}
\end{align}
Hence $c>0$ and $t=72d/c$. The determinant condition gives
$k^2\det T=-d$, so $\det T<0$ and $k$ is precisely one of
the two choices in \eqref{sgd:parameters2}. The scalar conditions
give precisely the two signs of $e$ stated there. The original
$\omega=e/k$ is unchanged, with sign $\epsilon\sigma$.

Conversely, start with parameters in $\mathcal D$. Positivity of $M$
follows from the exact factorization
\[
 M=LL^*,\qquad
 L=\begin{pmatrix}\sqrt u&0\\\bar m/\sqrt u&\sqrt d/\sqrt u\end{pmatrix},
 \qquad \det L=\sqrt d>0.
\]
The specified positive $t$ makes \eqref{sgd:trace-condition} zero.
Put $R=L^{-1}TL^{-*}$. This is Hermitian and
$\tr R=\tr(M^{-1}T)=0$, so its full form is
\[
 R=\begin{pmatrix}a&z\\\bar z&-a\end{pmatrix},\qquad a\in\mathbb R.
\]
The nonzero $m$ makes $T_{12}=-3m/2\ne0$, hence $R\ne0$.
Thus $\det R=-a^2-|z|^2<0$ and $\det T=d\det R<0$.
This proves that both signs of $k$ are well defined. Let $J_0=kR$.
It is Hermitian, has trace zero and determinant $-1$, and multiplying
the displayed matrix shows $J_0^2=I_2$.
For either sign and every column $z_0$, one has
\[
 z_0^*(I_2\pm J_0)z_0=\frac12\|(I_2\pm J_0)z_0\|^2\ge0.
\]
Also $\det(I_2\pm J_0)=0$ and $\tr(I_2\pm J_0)=2$,
so these matrices have rank one. Therefore
\[
 A_\pm=M\pm N=L(I_2\pm J_0)L^*
\]
are positive semidefinite of rank one, as claimed.

Their first diagonal entries are nonnegative. Suppose $h_s=0$ for
a sign $s\in\{1,-1\}$, where $A_s=M+sN$. Its determinant is
zero, so $|(A_s)_{12}|^2=0$. But
\[
 (A_s)_{12}=m(1-3sk/2),\quad
 h_s=u+sk(t/6-u).
\]
Since $m\ne0$, this forces $sk=2/3$ and then
$h_s=u/3+t/9>0$, a contradiction. Thus $h_\pm>0$.
Using $\det A_\pm=0$ gives
$(A_\pm)_{22}=|(A_\pm)_{21}|^2/h_\pm$; consequently
$A_\pm=a_\pm a_\pm^*$ with exactly the vectors in
\eqref{sgd:vectors}.

For those vectors, $v-iw=p_+a_+$ and $v+iw=p_-a_-$.
Their two outer products are $A_+,A_-$ because the phases have unit
modulus. Adding and subtracting the identities in
Lemma~\ref{sgd:realization-necessity} proves
\[
 vv^*+ww^*=M,\qquad i(vw^*-wv^*)=N.
\]
Equation~\eqref{sgd:scalar-conditions} holds by construction, so both
full scalar residuals vanish. Substitution of $N=kT$ and
$\omega k=e$ gives, with all coefficients retained,
\[
 \mathcal G=etE-2e(QM+MQ)
        -4e\left(\frac t4E-\frac12(QM+MQ)\right)=0.
\]
The exact identity $\gamma_0=-\tr\mathcal G$ then checks the
remaining upper Gauss entry. Thus every entry of all three original
residuals vanishes after applying \eqref{sgd:frame}.

For bijectivity, $u,v_0,y$ are recovered from $M$, $\epsilon$ is
the sign of $e$, and $\sigma$ the sign of $e/\omega$. The two
phase factors are recovered without any zero denominator by
\[
 p_+=\frac{(v-iw)_1}{\sqrt{h_+}},\qquad
 p_-=\frac{(v+iw)_1}{\sqrt{h_-}}.
\]
Indeed a column $c$ with $cc^*=a_\pm a_\pm^*$ and $h_\pm>0$
has $|c_1|=\sqrt{h_\pm}$, and the $(2,1)$ entry forces its second
component to be the same unit multiple of $(a_\pm)_2$.
These formulas prove both surjectivity and the exact inverse, without
setting any spatial phase to a selected value.
\end{proof}

\section{A short exact family with every equation checked}

\begin{proposition}\label{sgd:explicit-point}
Choose arbitrary $p_+,p_-\in U(1)$ and independent
$\epsilon,\sigma\in\{1,-1\}$. Put
\[
 M=\delta\begin{pmatrix}1&1\\1&2\end{pmatrix},\quad
 t=24\delta,\quad
 e=\epsilon\sqrt{2+3\delta},\quad
 k=\frac{2\sigma}{3\sqrt{13}},\quad
 \omega=\frac{3\sigma\sqrt{13}}2e,
\]
\[
 N=\frac{\sigma\delta}{\sqrt{13}}
        \begin{pmatrix}2&-1\\-1&-6\end{pmatrix},\qquad
 a_\pm=\sqrt\delta\begin{pmatrix}
 \sqrt{1\pm2\sigma/\sqrt{13}}\\[1mm]
 (1\mp\sigma/\sqrt{13})/\sqrt{1\pm2\sigma/\sqrt{13}}
 \end{pmatrix}.
\]
Define $v,w$ by \eqref{sgd:vectors}. For any retained unit frame
phase $\zeta$, these give original tuples by
\[
 C_0=B\bigl(ie\diag(1,0,-1)\bigr)B^*,\qquad
 Z_{1,0}=BX(v)B^*,\qquad Z_{2,0}=BX(w)B^*.
\]
Every such tuple satisfies \eqref{sgd:literal-H}--\eqref{sgd:literal-G}
at the positive original coupling $g^2=24\delta$.
\end{proposition}
\begin{proof}
The first diagonal factors are positive because $2/\sqrt{13}<1$.
For a sign $s$, the exact identity
\[
 (1-s\sigma/\sqrt{13})^2
  =(1+2s\sigma/\sqrt{13})(2-6s\sigma/\sqrt{13})
\]
follows by multiplying out and using $\sigma^2=1$. Hence
$a_\pm a_\pm^*=M\pm N$. The construction therefore has exactly
the displayed $M,N$ for every pair of phases. In particular
\[
 u=\delta,\quad v_0=2\delta,\quad m=\delta,\quad
 \tr M=3\delta,\quad d=\delta^2,\quad
 b=\tfrac12\tr N=-\frac{2\sigma\delta}{\sqrt{13}}.
\]
The scalar identities are $e^2=2+3\delta$ and
$2r\delta=1$; substitution in \eqref{sgd:scalar-H} and
\eqref{sgd:scalar-q} makes every entry zero.
For the full Gauss block, direct multiplication gives
\[
 QM+MQ=\delta\begin{pmatrix}2&3\\3&8\end{pmatrix},\qquad
 tE-2(QM+MQ)
   =\delta\begin{pmatrix}12&-6\\-6&-36\end{pmatrix}.
\]
The same matrix is $(4\omega/e)N$, since
$\omega/e=3\sigma\sqrt{13}/2$. Thus $\mathcal G=0$ entry
by entry, and the upper entry is explicitly
\[
 \gamma_0=4e\delta+20e\delta
       +8\left(\frac{3\sigma\sqrt{13}}2e\right)
             \left(-\frac{2\sigma\delta}{\sqrt{13}}\right)=0.
\]
Finally \eqref{sgd:frame} proves all three literal equations exactly.
As an additional check on the construction, its determinant data are
\[
 T=\frac{3\delta}{2}\begin{pmatrix}2&-1\\-1&-6\end{pmatrix},\quad
 \det T=-\frac{117}{4}\delta^2,\quad k^2=\frac4{117},
\]
agreeing with the general parameterization. Both signs of the original
frequency are present because $\operatorname{sign}\omega=\epsilon\sigma$.
\end{proof}

\begin{corollary}[Every positive original coupling occurs]
\label{sgd:every-coupling}
For every specified original $\tau=g^2>0$, the literal two-cross
scalar--Gauss domain contains tuples at $t=\tau$, with either original
frequency sign and both temporal signs. It contains continuous phase
families at each such coupling.
\end{corollary}
\begin{proof}
For $\rho>1/\sqrt2$ choose
\[
 M_\rho=\delta\begin{pmatrix}\rho&1\\1&2\rho\end{pmatrix}.
\]
Its parameters obey $u=\delta\rho>0$,
$d=\delta^2(2\rho^2-1)>0$, $c=3\delta\rho>0$ and $m=\delta$.
Thus Theorem~\ref{sgd:parametrization} gives the complete solution
family at
\[
 t=\frac{72\delta^2(2\rho^2-1)}{3\delta\rho}
   =\frac{24\delta(2\rho^2-1)}\rho.
\]
For a given $\tau>0$, set
\[
 \rho=\frac{\tau+\sqrt{\tau^2+4608\delta^2}}{96\delta}.
\]
Because $\sqrt{4608}\delta=48\sqrt2\delta$, this is strictly
greater than $1/\sqrt2$. Squaring the displayed defining equality
gives $48\delta\rho^2-\tau\rho-24\delta=0$, which is precisely
$t=\tau$. No limiting passage or change of coupling is used.
The independent choices of $\epsilon,\sigma,p_+,p_-$ supply the
stated signs and phases; the inverse in Theorem~\ref{sgd:parametrization}
ensures different phase pairs yield different classified tuples.
\end{proof}

\section{Exact scope and reproducibility}

The results concern the simultaneous zero set of both scalar equations
and Gauss. The construction never imposes either spatial equation and
does not claim a full rotating-ansatz solution. It proves that the
previously specified scalar--Gauss domain of the two-cross quantitative
residual certificate is nonempty, in fact at every original positive
coupling. This is an algebraic statement about the retained classical
residuals and makes no quantum assertion.

The companion \texttt{check\_scalar\_gauss\_domain.py} is a finite exact
symbolic replay. It checks the full scalar and Gauss matrix formulas with
eight independent real spatial coordinates, the determinant/trace
identities, the explicit factors and both signs, and the original
unitary frame maps. The rank, positivity, parameter domains, completeness
and inverse-map proofs are given in this document; they are not inferred
from sampled values or from a numerical residual tolerance.

\end{document}
